\section{The local expansion: addition theorem, exact moments, sub-boxes}
\label{sec:exp}

This section expands the volume form \eqref{eq:vol:main} in the size of the cell. The
mechanism is the spherical addition theorem: it separates the singular dependence on the
centre separation $\mathbf{R}$, where the kernel is elementary, from the regular
dependence on the difference variable $\boldsymbol{\delta}$, where the kernel is a
homogeneous polynomial whose integral against the triangle weight is an exact rational in
the edge lengths. The result is a two-index expansion --- a multipole order $l$ and a
frequency order $n$ --- whose geometry data are frequency independent and computed once
per cell shape, and whose per-offset work is a sum of elementary functions.

\subsection{The addition theorem, and the sign the usual statement hides}
\label{sec:exp:add}

Write $r = |\mathbf{R}+\boldsymbol{\delta}|$, $\hat{\mathbf{R}} = \mathbf{R}/|\mathbf{R}|$,
$\hat{\boldsymbol{\delta}} = \boldsymbol{\delta}/|\boldsymbol{\delta}|$.

\begin{proposition}[the form actually needed]
\label{prop:exp:add}
For $|\boldsymbol{\delta}| < |\mathbf{R}|$,
\begin{align}
\frac{e^{\ii k r}}{r}
&= \ii k \sum_{l\ge 0} (2l+1)\,(-1)^l\,
   j_l(k|\boldsymbol{\delta}|)\,h^{(1)}_l(k|\mathbf{R}|)\,
   P_l\bigl(\hat{\mathbf{R}}\cdot\hat{\boldsymbol{\delta}}\bigr)
\label{eq:exp:addP}\\
&= 4\pi\,\ii k \sum_{l\ge 0}\sum_{m=-l}^{l} (-1)^l\,
   h^{(1)}_l(k|\mathbf{R}|)\,Y_{lm}(\hat{\mathbf{R}})\,
   j_l(k|\boldsymbol{\delta}|)\,Y_{lm}(\hat{\boldsymbol{\delta}}) .
\label{eq:exp:addY}
\end{align}
\end{proposition}

\begin{remark}[the convention, stated once]
\label{rem:exp:sign}
The classical theorem is usually written for $\mathbf{r} = \mathbf{r}_1 - \mathbf{r}_2$
with $P_l(\cos\gamma)$ and $\gamma$ the angle between $\mathbf{r}_1$ and $\mathbf{r}_2$.
Here $\mathbf{r} = \mathbf{R} - (-\boldsymbol{\delta})$, so $\mathbf{r}_2 =
-\boldsymbol{\delta}$ and $\cos\gamma = -\hat{\mathbf{R}}\cdot\hat{\boldsymbol{\delta}}$;
since $P_l(-t) = (-1)^l P_l(t)$, the form written in terms of
$P_l(\hat{\mathbf{R}}\cdot\hat{\boldsymbol{\delta}})$ --- the one whose angular factor is
the real solid harmonic of $\boldsymbol{\delta}$ itself, which is what the moments need
--- \emph{carries an explicit $(-1)^l$}. Omitting it is a first-order error, not a
subtlety: at $\mathbf{R} = (1,2,3)$, $\boldsymbol{\delta} = (0.3,-0.2,0.1)$, $f=1$, in
$300$-bit arithmetic with $L = 40$, \eqref{eq:exp:addY} reproduces $e^{\ii kr}/r$ to
$1.06\times10^{-41}$ with the factor and to a relative error of $6.66\times10^{-1}$
without it. The constant $4\pi\ii k$ is therefore verified to $1.1\times10^{-41}$ at
$f=1$ and $1.1\times10^{-40}$ at $f = 1+0.1\ii$. For the whole difference box the factor
turns out to be moot, because Section~\ref{sec:exp:parity} shows that every surviving $l$
is even; \emph{on a sub-box it is mandatory}, because subdivision destroys that parity
(Section~\ref{sec:exp:sub}).
\end{remark}

Multiplying \eqref{eq:exp:addY} by $1/(4\pi f^2)$ gives the kernel,
\begin{equation}
g(|\mathbf{R}+\boldsymbol{\delta}|)
= \frac{\ii k}{f^2}\sum_{l,m}(-1)^l\,h_l(k|\mathbf{R}|)\,Y_{lm}(\hat{\mathbf{R}})\,
  j_l(k|\boldsymbol{\delta}|)\,Y_{lm}(\hat{\boldsymbol{\delta}}) ,
\label{eq:exp:g}
\end{equation}
absolutely and uniformly convergent on compact subsets of
$|\boldsymbol{\delta}| < |\mathbf{R}|$, so it may be integrated against $w$ and
differentiated in $\mathbf{R}$ term by term. Throughout, $h_l$ means $h_l^{(1)} = j_l +
\ii y_l$.

\subsection{Real solid harmonics as integer polynomials}
\label{sec:exp:solid}

Take real harmonics without the Condon--Shortley phase,
\begin{equation}
Y_{l0} = N_{l0}P_l(\cos\theta),\quad
Y_{lm} = \sqrt2\,N_{lm}P_l^m(\cos\theta)\cos m\varphi,\quad
Y_{l,-m} = \sqrt2\,N_{lm}P_l^m(\cos\theta)\sin m\varphi,
\label{eq:exp:Y}
\end{equation}
$N_{lm} = \sqrt{(2l+1)(l-m)!/(4\pi(l+m)!)}$, so that $\sum_m Y_{lm}(\mathbf{a})
Y_{lm}(\mathbf{b}) = \frac{2l+1}{4\pi}P_l(\mathbf{a}\cdot\mathbf{b})$. The regular solid
harmonic $|\boldsymbol{\delta}|^l Y_{lm}(\hat{\boldsymbol{\delta}})$ is then an
\emph{integer} polynomial in $(\delta_1,\delta_2,\delta_3)$ times a single square root:
\begin{equation}
|\boldsymbol{\delta}|^{l}Y_{lm}(\hat{\boldsymbol{\delta}})
= \nu_{lm}\;C_l^m\bigl(\delta_3, |\boldsymbol{\delta}|^2\bigr)\;A_m(\delta_1,\delta_2),
\qquad
\nu_{lm} = \sqrt{\frac{(2l+1)(2-\delta_{m0})}{4\pi\,(l-m)!\,(l+m)!}},
\label{eq:exp:solid}
\end{equation}
with $A_m = \operatorname{Re}(\delta_1+\ii\delta_2)^{m}$ for the cosine family and
$A_m = \operatorname{Im}(\delta_1+\ii\delta_2)^{m}$ for the sine family, and with
$C_l^m$ generated by the integer recurrence
\begin{equation}
C_m^m = (2m-1)!!,\qquad C_{m+1}^m = (2m+1)!!\,\delta_3,\qquad
C_l^m = (2l-1)\,\delta_3\,C_{l-1}^m - (l+m-1)(l-m-1)\,|\boldsymbol{\delta}|^2\,C_{l-2}^m .
\label{eq:exp:rec}
\end{equation}
Note the $|\boldsymbol{\delta}|^2$ and not $\delta_1^2+\delta_2^2$: $C_l^m$ is
$(l-m)!$ times $|\boldsymbol{\delta}|^l P_l^m$ stripped of the factor
$(\delta_1+\ii\delta_2)^m$, and it is exactly the factorial $(l-m)!$ that makes
\eqref{eq:exp:rec} integral, cancelling afterwards against $\nu_{lm}$. Stored in the
parity-compact array $P[a+1,b+1] = $ the coefficient of
$\delta_1^{2a+p_1}\delta_2^{2b+p_2}\delta_3^{\deg-\cdots}$, multiplication by $\delta_3$
is free (the degree label carries it) and multiplication by $|\boldsymbol{\delta}|^2$ is
$\mathrm{dst}[a,b] = \mathrm{src}[a,b]+\mathrm{src}[a-1,b]+\mathrm{src}[a,b-1]$. Checked
against the normalised associated-Legendre recursion at a generic point for every $(l,m)$
with $l \le 20$ at $300$ bits: maximum relative difference $2.8\times10^{-87}$.

The recurrence \eqref{eq:exp:rec} alternates in sign and the geometry contraction that
follows cancels badly --- $\sum|\text{term}|/|\text{result}| = 2.1\times10^{5}$ at
$L = 20$ and $5.8\times10^{8}$ at $L = 40$ for the cubes, $1.1\times10^{7}$ for the
slender cell, i.e.\ up to $8.8$ decimal digits. That is why the geometry table is built in
exact rational arithmetic once per cell shape (\code{Rational} of \code{BigInt} edge lengths,
\code{BigInt} polynomial coefficients) and rounded to \code{Float64} only at the end. The
cost is a one-off: $14.65$~s at $L = 40$, $n_{\max} = 12$.

\subsection{\texorpdfstring{The regular factor as a power series in $k$}{The regular factor as a power series in k}}
\label{sec:exp:kser}

$j_l$ is entire and even in its argument up to the prefactor $z^l$:
\begin{equation}
j_l(k|\boldsymbol{\delta}|)\,Y_{lm}(\hat{\boldsymbol{\delta}})
= \sum_{n\ge 0} c_{ln}(k)\;
  |\boldsymbol{\delta}|^{2n}\bigl[\,|\boldsymbol{\delta}|^{l}
  Y_{lm}(\hat{\boldsymbol{\delta}})\,\bigr],
\qquad
c_{ln}(k) = \frac{(-1)^n\,k^{\,l+2n}}{2^n\,n!\,(2l+2n+1)!!},
\label{eq:exp:cln}
\end{equation}
which follows from $j_l(z) = \frac{z^l}{(2l+1)!!}\sum_n \frac{(-z^2/2)^n}{n!\prod_{q=1}^{n}
(2l+2q+1)}$ and $(2l+1)!!\prod_{q=1}^n(2l+2q+1) = (2l+2n+1)!!$. Every summand of
\eqref{eq:exp:cln} is a homogeneous polynomial in $\boldsymbol{\delta}$ of degree
$l + 2n$ with rational coefficients, so its integral against $w$ is exact. That is the
whole point: \emph{$k$ enters only through the scalars $c_{ln}$ and $h_l(k|\mathbf{R}|)$;
the geometry is frequency independent.}

\subsection{Where the derivatives go}
\label{sec:exp:der}

Because $g$ depends on $\mathbf{R}$ and $\boldsymbol{\delta}$ only through their sum,
$\partial/\partial R_a$ and $\partial/\partial\delta_a$ act identically on
$g(\mathbf{R}+\boldsymbol{\delta})$, and the two derivatives of \eqref{eq:vol:main} may be
put on either side. Both organizations are used below, for different purposes.

\paragraph{Derivatives on the regular side (the evaluation).} Acting on the polynomial
\eqref{eq:exp:cln} the derivatives only re-index the moment product, since
$\partial_a^2 \boldsymbol{\delta}^{\boldsymbol{\alpha}} = \alpha_a(\alpha_a-1)
\boldsymbol{\delta}^{\boldsymbol{\alpha}-2\mathbf{e}_a}$ and
$\partial_a\partial_b\boldsymbol{\delta}^{\boldsymbol{\alpha}} = \alpha_a\alpha_b
\boldsymbol{\delta}^{\boldsymbol{\alpha}-\mathbf{e}_a-\mathbf{e}_b}$. Define the three
exact-rational geometry tables
\begin{equation}
\begin{aligned}
A^{\mathrm{d}}[a,n,lm] &= \frac{1}{V_{\mathrm{t}}}\int_D w\;
  \partial_a^2\bigl[\,|\boldsymbol{\delta}|^{2n+l}Y_{lm}\,\bigr]\,\dd\boldsymbol{\delta},
&\quad a &= 1,2,3,\\
B^{\mathrm{d}}[n,lm] &= \frac{1}{V_{\mathrm{t}}}\int_D w\;
  |\boldsymbol{\delta}|^{2n+l}Y_{lm}\,\dd\boldsymbol{\delta},\\
A^{\mathrm{o}}[ab,n,lm] &= \frac{1}{V_{\mathrm{t}}}\int_D w\;
  \partial_a\partial_b\bigl[\,|\boldsymbol{\delta}|^{2n+l}Y_{lm}\,\bigr]\,
  \dd\boldsymbol{\delta}, &\quad ab &= 12, 13, 23,
\end{aligned}
\label{eq:exp:tabs}
\end{equation}
each an exact rational in $s_1,s_2,s_3$ times $\nu_{lm}$, by \eqref{eq:vol:mu}. Contract
them with the frequency once per $f$,
\begin{equation}
\widetilde{Q}^{\,aa}_{lm}(k) = \sum_{n} c_{ln}(k)\bigl(A^{\mathrm{d}}[a,n,lm]
  + k^2 B^{\mathrm{d}}[n,lm]\bigr),
\qquad
\widetilde{Q}^{\,ab}_{lm}(k) = \sum_{n} c_{ln}(k)\,A^{\mathrm{o}}[ab,n,lm],
\label{eq:exp:Q}
\end{equation}
and then \emph{per offset} only the elementary sum remains,
\begin{equation}
\boxed{\;
T_{ab}(\mathbf{R}) \;=\; \frac{\ii k}{f^2}
\sum_{l\,\le\,L}\ \sum_{m=-l}^{l} (-1)^l\,
h_l(k|\mathbf{R}|)\,Y_{lm}(\hat{\mathbf{R}})\;\widetilde{Q}^{\,ab}_{lm}(k).\;}
\label{eq:exp:eval}
\end{equation}
The frequency contraction \eqref{eq:exp:Q} costs $31$~ms once per frequency; the per-offset
cost of \eqref{eq:exp:eval} is $391$, $1250$, $2731$ and $4972$~ns at $L = 10, 20, 30,
40$, i.e.\ $t(L) \approx 3.1L^2 + 80$~ns with no allocation.

\paragraph{Derivatives on the singular side (the bound).} The same object is
\begin{equation}
T_{ab}(\mathbf{R}) = \frac{\ii k}{f^2 V_{\mathrm{t}}}\sum_{l,m}(-1)^l
\Bigl[(\partial_a\partial_b + \delta_{ab}k^2)
\bigl(h_l(k|\mathbf{R}|)Y_{lm}(\hat{\mathbf{R}})\bigr)\Bigr]\,\mathcal{I}_{lm},
\qquad
\mathcal{I}_{lm} = \int_D w\, j_l(k|\boldsymbol{\delta}|)Y_{lm}(\hat{\boldsymbol{\delta}})\,
\dd\boldsymbol{\delta},
\label{eq:exp:sing}
\end{equation}
and it is this form that Section~\ref{sec:bnd} bounds, because the $\boldsymbol{\delta}$
dependence has been reduced to a single non-negative radial moment while the derivatives
sit on functions of $\mathbf{R}$ alone, whose growth in $l$ is known in closed form.
Note that the two organizations select different $(l,m)$: in \eqref{eq:exp:sing} the
surviving classes are only those of the diagonal (Section~\ref{sec:exp:parity}), the six
tensor entries being produced by the derivative; in \eqref{eq:exp:eval} the derivative has
already been absorbed into the tables and four classes survive.

\subsection{What the triangle weight kills}
\label{sec:exp:parity}

Each monomial $\delta_1^{i}\delta_2^{j}\delta_3^{p}$ of
$|\boldsymbol{\delta}|^{2n}\,|\boldsymbol{\delta}|^{l}Y_{lm}$ has a parity triple fixed by
$(l,m)$ and by the family:
\begin{equation}
\operatorname{par}(\delta_1) = (m-t)\bmod 2,\qquad
\operatorname{par}(\delta_2) = t \bmod 2,\qquad
\operatorname{par}(\delta_3) = (l-m)\bmod 2,
\label{eq:exp:par}
\end{equation}
with $t$ even for the cosine family and odd for the sine family. Since $w$ is even in
every coordinate, $\int_D w\,\boldsymbol{\delta}^{\boldsymbol{\alpha}} \neq 0$ only when
every $\alpha_i$ is even, and $\partial_a\partial_b$ flips the parity of $\alpha_a$ and
$\alpha_b$. Exactly four classes survive:
\begin{center}\small
\begin{tabular}{@{}ll@{}}
\toprule
entry & surviving class\\
\midrule
$11$, $22$, $33$ & cosine family, $m$ even $\ge 0$, $l-m$ even $\;\Rightarrow\; l$ even\\
$12$ & sine family, $m$ even $\ge 2$, $l-m$ even $\;\Rightarrow\; l$ even\\
$13$ & cosine family, $m$ odd, $l-m$ odd $\;\Rightarrow\; l$ even\\
$23$ & sine family, $m$ odd, $l-m$ odd $\;\Rightarrow\; l$ even\\
\bottomrule
\end{tabular}
\end{center}
\noindent
\emph{Every surviving $l$ is even}, which is why the $(-1)^l$ of
Remark~\ref{rem:exp:sign} never appears for the whole box. The stored count is $144$,
$484$, $1024$, $1764$ numbers per $n$ at $L = 10, 20, 30, 40$ against $121$, $441$,
$961$, $1681$ for one full $(l,m)$ table --- the four classes together cost about what a
single unpruned table would, while carrying all three diagonal derivatives and both
off-diagonal families.

\begin{remark}[the cube loses $l=2$ as well]
\label{rem:exp:l2}
For a \emph{cubic} cell the entire $l=2$ shell vanishes identically, so the surviving
shells are $l \in \{0,4,6,8,\dots\}$. In the organization \eqref{eq:exp:sing} this is
immediate: $|\boldsymbol{\delta}|^2Y_{20} \propto 2\delta_3^2-\delta_1^2-\delta_2^2$ and
$|\boldsymbol{\delta}|^2Y_{22} \propto \delta_1^2-\delta_2^2$, and
$\int_D w\,\delta_i^2 = \mu_2(s_i)\prod_{j\neq i}\mu_0(s_j)$, which is the same number for
$i=1,2,3$ when $s_1=s_2=s_3$; the two combinations therefore cancel exactly. Measured:
the shells with $\max_m|\mathcal{I}_{lm}|$ above $10^{-40}$ of the largest are
$\{0,4,6,8,\dots\}$ for $(1/32)^3$ and $(1/4)^3$ and $\{0,2,4,6,\dots\}$ for
$(1/32,1/32,1/512)$ (script \code{chkRem.jl}). This is a genuine saving on the shapes
Gila uses most and is not exploited by the term-count tables of
Section~\ref{sec:bnd}, which are stated for the generic shape.
\end{remark}

\subsection{Internal identities of the geometry table}
\label{sec:exp:ident}

Two identities test the tables \eqref{eq:exp:tabs} without any reference value. First,
harmonicity: $|\boldsymbol{\delta}|^lY_{lm}$ is harmonic, so
$\sum_a A^{\mathrm{d}}[a,0,lm] = 0$, and in rational arithmetic this is exactly zero.
Second, the Helmholtz identity term by term: for $H_{lm} =
|\boldsymbol{\delta}|^lY_{lm}$,
\begin{equation}
\begin{gathered}
\nabla^2\bigl[\,|\boldsymbol{\delta}|^{2n}H_{lm}\,\bigr]
= 2n\,(2n+2l+1)\,|\boldsymbol{\delta}|^{2n-2}H_{lm},\\
\text{hence}\qquad
\sum_a A^{\mathrm{d}}[a,n,lm] = 2n(2n+2l+1)\,B^{\mathrm{d}}[n-1,lm].
\end{gathered}
\label{eq:exp:helm}
\end{equation}
On the exact-rational table with $l \le 20$, $1 \le n \le 6$, the maximum relative
violation is $1.73\times10^{-89}$ at $300$ bits of output precision, i.e.\ the identity
holds exactly and the printing is the only error. Independently, the constant $4\pi\ii k$
and the $(-1)^l$ of \eqref{eq:exp:addY} are verified to $1.7\times10^{-89}$ on the full
assembled identity.

\subsection{\texorpdfstring{Octant sub-boxes: affine weights and the moments $\nu_n$}{Octant sub-boxes: affine weights and the moments nu\_n}}
\label{sec:exp:sub}

The expansion \eqref{eq:exp:eval} converges like $\rho^l$ with $\rho =
r_{\mathrm{d}}/|\mathbf{R}|$, $r_{\mathrm{d}} = |\mathbf{s}|$, and at the nearest
non-touching offsets $\rho$ is close to one. Subdividing $D$ shrinks it. Split each axis
into pieces, so that $D$ is the disjoint union of boxes $D_j$ with centres
$\mathbf{c}_j$ and half-widths $\mathbf{h}_j$; with $\boldsymbol{\delta}' =
\boldsymbol{\delta}-\mathbf{c}_j$,
\begin{equation}
T_{ab}(\mathbf{R}) = \sum_j T^{(j)}_{ab}(\mathbf{R}+\mathbf{c}_j),
\qquad
T^{(j)}_{ab}(\mathbf{V}) = \frac{1}{V_{\mathrm{t}}}\int_{D_j}
w(\mathbf{c}_j+\boldsymbol{\delta}')\,
\bigl[(\partial_a\partial_b+\delta_{ab}k^2)g\bigr](\mathbf{V}+\boldsymbol{\delta}')\,
\dd\boldsymbol{\delta}' ,
\label{eq:exp:split}
\end{equation}
each piece being the same object with local centre $\mathbf{V} = \mathbf{R}+\mathbf{c}_j$,
local radius $r_j = |\mathbf{h}_j|$ and local ratio $\rho_j = r_j/|\mathbf{V}|$.

\begin{lemma}[no piece may straddle the origin]
\label{lem:exp:straddle}
On a piece $[a,b]$ of axis $i$ with $a \ge 0$ the factor $s_i - |\delta_i|$ is
$\alpha_i + \beta_i t$ with $t = \delta'_i \in [-h_i,h_i]$,
$\alpha_i = s_i - |c_i|$ and $\beta_i = -\operatorname{sign}(c_i)$; likewise for
$b \le 0$. A piece containing $0$ in its interior carries the kink of $|\delta_i|$ and has
no affine form, and using one anyway is not a small error: for $m = 3$ equal pieces the
middle piece $[-s/3,s/3]$ has exact zeroth moment $5s^2/9$ against $2s^2/3$ for the best
affine model, a $20\,\%$ error on that piece.
\end{lemma}

Hence $0$ must be a split point on every axis. With $m_i$ equal pieces that is automatic
for even $m_i$ and forces an extra split for odd $m_i$, which then has $m_i+1$ pieces:
\emph{$m = 3$ costs exactly what $m = 4$ costs and converges more slowly}, its outer
pieces having half-width $s/3$ against $s/4$, so odd splits are dominated and never used.
The minimal legal split is the eight octants, $m = (2,2,2)$. On a piece the moments are
again exact rationals,
\begin{equation}
\nu_n(\alpha,\beta,h) = \int_{-h}^{h}(\alpha+\beta t)\,t^n\,\dd t
= \begin{cases} \dfrac{2\alpha h^{\,n+1}}{n+1}, & n \text{ even},\\[4pt]
                \dfrac{2\beta h^{\,n+2}}{n+2}, & n \text{ odd},\end{cases}
\qquad
\int_{D_j} w\,\boldsymbol{\delta}'^{\boldsymbol{\gamma}}\,\dd\boldsymbol{\delta}'
= \prod_i \nu_{\gamma_i}(\alpha_i,\beta_i,h_i),
\label{eq:exp:nu}
\end{equation}
and only three one-dimensional tables are needed, $\nu^{(0)}_n = \nu_n$,
$\nu^{(1)}_n = n\,\nu_{n-1}$, $\nu^{(2)}_n = n(n-1)\nu_{n-2}$, exactly as for the whole
box. Two exact consistency checks hold with zero residual in
\code{Rational} of \code{BigInt}: $\sum_{\text{pieces}}\nu_0 = \mu_0(s)$ for
$m = 1,2,3,4,5,8$, and $\sum_j\sum_q \binom{n}{q}c_j^{\,n-q}\nu_q = \mu_n(s)$ for
$n \le 8$, $m = 2,3,4,6$.

Subdivision destroys the parity argument of Section~\ref{sec:exp:parity}: the affine
weight $\alpha+\beta t$ on an off-centre piece is not even, so \emph{every} $(l,m)$ with
$0 \le l \le L$ contributes to every one of the six entries (seven kinds $\times$
$(n_{\max}+1)$ $\times$ $(L+1)^2$ numbers per sub-box type, and about $7$ complex
multiplies per $(l,m)$ against roughly $1.7$ for the whole box), and \emph{odd $l$
survives, so the $(-1)^l$ becomes mandatory}. Omitting it is a first-order error: at
$(1/32)^3$, $m = (2,2,2)$, $L = 24$, $f = 1$, the relative discrepancy against the
unsubdivided value is $3.65\times10^{-2}$ at $\mathbf{n} = (4,0,0)$,
$5.63\times10^{-2}$ at $(3,1,0)$, $6.45\times10^{-3}$ at $(6,6,6)$ and
$7.58\times10^{-3}$ at $(5,1,2)$, against $1.51\times10^{-13}$,
$2.63\times10^{-11}$, $2.85\times10^{-16}$ and $2.59\times10^{-16}$ with the factor
restored --- and the first two residuals are the unsubdivided value's own $l$-truncation
at $L = 24$, not the subdivided one's.

\subsection{The reflection rule: one octant of tables}
\label{sec:exp:refl}

\begin{proposition}[reflection]
\label{prop:exp:refl}
Let $\sigma \in \{\pm1\}^3$ act on coordinates. Then
\begin{equation}
T^{(\sigma j)}_{ab}(\mathbf{V}) = \sigma_a\sigma_b\,T^{(j)}_{ab}(\sigma\mathbf{V}),
\qquad\text{so}\qquad
T_{ab}(\mathbf{R}) = \sum_{t \in \text{one octant}}\ \sum_{\sigma\in\{\pm1\}^3}
\sigma_a\sigma_b\,T^{(t)}_{ab}(\sigma\mathbf{R}+\mathbf{c}_t).
\label{eq:exp:refl}
\end{equation}
\end{proposition}

\begin{proof}
$w$ is even, so $w(\sigma\mathbf{x}) = w(\mathbf{x})$, and
$G_{ab}(\mathbf{x}) = [(\partial_a\partial_b+\delta_{ab}k^2)g](\mathbf{x})$ satisfies
$G_{ab}(\sigma\mathbf{y}) = \sigma_a\sigma_b G_{ab}(\mathbf{y})$ because $g$ depends only
on $|\mathbf{y}|$ and the two derivatives carry one sign each. Substituting
$\boldsymbol{\delta}' = \sigma\mathbf{u}$ in \eqref{eq:exp:split} gives the first
statement; the sub-box set is closed under $\sigma$ with
$\mathbf{c}_{\sigma j} = \sigma\mathbf{c}_j$, which gives the second.
\end{proof}

Diagonal entries take no sign and the off-diagonal $(a,b)$ flips whenever exactly one of
the axes $a,b$ is flipped --- the same rule the assembled tensor obeys under a lattice
reflection. Consistently, the mirrored moments are
$\nu_n(\alpha,-\beta,h) = (-1)^n\nu_n(\alpha,\beta,h)$: the mirrored table is the original
with every odd-$\gamma_i$ monomial negated. So only
$\prod_i(\text{pieces per axis on the positive side})$ geometry tables are ever built ---
one for the octant split. Verified against tables built independently for all eight
octants at $(1/32)^3$, $m = (2,2,2)$, $L = 16$: $1.70\times10^{-16}$ at
$\mathbf{n} = (3,1,0)$, $1.35\times10^{-16}$ at $(2,1,1)$, $9.61\times10^{-17}$ at
$(5,1,2)$. Per sub-box type and sign pattern, with $\mathbf{V} = \sigma\mathbf{R} +
\mathbf{c}_t$,
\begin{equation}
T^{(t)}_{ab}(\mathbf{V}) = \frac{\ii k}{f^2}\sum_{l=0}^{L_{t,\sigma}}\sum_{m=-l}^{l}
(-1)^l h_l(k|\mathbf{V}|)\,Y_{lm}(\hat{\mathbf{V}})\,\widetilde{Q}^{\,ab}_{lm}(t,k),
\label{eq:exp:subeval}
\end{equation}
with $\widetilde{Q}$ built from \eqref{eq:exp:nu} exactly as \eqref{eq:exp:Q} is built
from \eqref{eq:vol:mu}, and with the $(-1)^l$ folded into $\widetilde{Q}$ once.

\subsection{Real and imaginary parts as separate real sums}
\label{sec:exp:reim}

For real $f$ the coefficients $c_{ln}(k)$ of \eqref{eq:exp:cln} are real, the geometry
tables \eqref{eq:exp:tabs} are real, and hence $\widetilde{Q}^{\,ab}_{lm}$ is real. With
$h_l = j_l + \ii y_l$ and $\ii k/f^2 = 2\pi\ii/f$, \eqref{eq:exp:eval} splits into two
independent real sums,
\begin{equation}
\operatorname{Re}T_{ab} = -\frac{k}{f^2}\sum_{l,m} y_l(k|\mathbf{R}|)\,
Y_{lm}(\hat{\mathbf{R}})\,\widetilde{Q}^{\,ab}_{lm},
\qquad
\operatorname{Im}T_{ab} = +\frac{k}{f^2}\sum_{l,m} j_l(k|\mathbf{R}|)\,
Y_{lm}(\hat{\mathbf{R}})\,\widetilde{Q}^{\,ab}_{lm}.
\label{eq:exp:reim}
\end{equation}
This is a structural advantage over any organization that forms one complex sum, and it is
worth stating because it is measurable. The radiative (imaginary) part of a far entry is
small compared with the reactive part: when $|\operatorname{Im}T_{ab}|/\max_{ab}|T_{ab}|$
falls below about $10^{-3}$ --- which happens routinely at low frequency and on slender
cells --- a single complex sum delivers the imaginary part only to
$\varepsilon\cdot|\operatorname{Re}T|$ in absolute terms, i.e.\ to a \emph{relative}
accuracy of $10^{-16}/(|\operatorname{Im}T|/\max|T|)$. Measured on a Cartesian-Taylor
route that does form one complex sum: relative imaginary errors of
$1.33\times10^{-13}$ and $1.31\times10^{-13}$ on the $\lambda/32$ cube at $f = 1+0.1\ii$
where the ratio is $3.6\times10^{-3}$ and $2.0\times10^{-3}$,
$2.37\times10^{-12}$ on the slender cell at $f = 0.37$ where it is
$1.6\times10^{-6}$, and $3.99\times10^{-12}$ at ratio $6.0\times10^{-5}$. The split
\eqref{eq:exp:reim} removes that coupling; it costs nothing, and it makes the
$y_l$ and $j_l$ recurrences separately responsible for their own part. The price is that
$j_l(k|\mathbf{R}|)$ must then be obtained to \emph{relative} accuracy for $l$ above
$k|\mathbf{R}|$, where the upward recurrence for $h_l$ destroys it (the error of
$\operatorname{Re}h_l$ from the upward recurrence is $1.3\times10^{13}$ relative at
$z = 0.3$, $l = 10$ and $1.4\times10^{24}$ at $z = 10$, $l = 40$); Miller's downward
recurrence returns $j_l$ to $6.8\times10^{-15}$ at worst over
$l \le 40$, $|z| \in [0.3,300]$, at a flat cost of $850$~ns.
